% Generated draft: requires mathematical and source review.
\documentclass[11pt]{article}
\usepackage[margin=25mm]{geometry}
\usepackage{fontspec}
\setmainfont{Latin Modern Roman}
\usepackage{amsmath,amssymb,mathtools,unicode-math}
\setmathfont{Latin Modern Math}
\usepackage{enumitem,needspace,xurl,hyperref}
\setlength{\parindent}{0pt}
\setlength{\parskip}{5pt}
\setlength{\emergencystretch}{3em}
\newcommand{\N}{\mathbb{N}}
\newcommand{\Z}{\mathbb{Z}}
\newcommand{\Q}{\mathbb{Q}}
\newcommand{\R}{\mathbb{R}}
\newcommand{\C}{\mathbb{C}}
\newcommand{\F}{\mathbb{F}}
\newcommand{\A}{\mathbb{A}}
\newcommand{\PP}{\mathbb{P}}
\newcommand{\E}{\mathbb{E}}
\newcommand{\eps}{\varepsilon}
\newcommand{\dd}{\,\mathrm{d}}
\begin{document}
\section*{UnsolvedMath ID: 20001639}
\textbf{Source problem code:} catalogue entry
\subsection{Definitions and mathematical statement}
Let \(F_n\) be a free group of rank \(n\), and \(\operatorname{Out}(F_n)=\operatorname{Aut}(F_n)/\operatorname{Inn}(F_n)\). An iwip is fully irreducible: no positive power preserves the conjugacy class of a nontrivial proper free factor. Polynomial growth refers to polynomial growth of the lengths of iterated conjugacy classes in a fixed free basis. For an isometry \(\gamma\) of a graph with its edge metric, write \(\tau(\gamma)=\lim_{k\to\infty}d(x,\gamma^kx)/k\); ellipticity means bounded orbits, a stronger general condition than merely \(\tau=0\).

The workshop asks for a \(\delta\)-hyperbolic graph with an \(\operatorname{Out}(F_n)\)-action such that every iwip has positive translation length; polynomially growing elements are elliptic; the construction is elegant, beautiful, useful, or truthful; local structure supports induction; and the complex and links are highly connected. The last three requirements have no numerical or formal specification. Connectivity and links refer naturally to a simplicial complex with the graph as one-skeleton, since higher connectivity of a graph alone would be a different demand. The source's parenthetical equality between ellipticity and zero translation length is therefore recorded as an informal shorthand rather than a general equivalence.

The supplied aid reports, for \(n\ge3\), a partial realization by the free-factor complex \(\mathcal F_n\), whose vertices are conjugacy classes of nontrivial proper free factors and whose simplices are chains of inclusions up to conjugacy. It reports hyperbolicity of its one-skeleton, \(\tau(\gamma)>0\) for every iwip, a finite vertex orbit for every polynomially growing element, and \((n-3)\)-connectivity of the complex. It further reports, for every corank-one free factor \(A\) of rank \(n-1\), a natural simplicial identification
\[
 \operatorname{Lk}_{\mathcal F_n}([A])\cong\mathcal F(A)\cong\mathcal F_{n-1}.
\]
This is the conjugacy-class inclusion complex convention, including the rank-two link convention; other modified low-rank graphs are not silently substituted. These assertions do not specify all links or settle every qualitative workshop requirement.

A precise working version of the aid's two-arrow transfer test is: for equivariant Lipschitz maps \(Z\to X\to Y\), positive stable lengths in \(Y\) force positive stable lengths in \(X\), while bounded orbits in \(Z\) force bounded orbits in \(X\). In particular an equivariant \(L\)-Lipschitz map \(X\to Y\) has \(\tau_Y(\gamma)\le L\tau_X(\gamma)\). The aid applies this to report, for even \(n=2g\ge4\), that no \(\operatorname{Out}(F_n)\)-equivariant Lipschitz map runs from the co-surface graph \(CS_{2g}\) to \(\mathcal F_{2g}\): geometric iwips are the stated dynamical obstruction. The source does not define the co-surface graph or give the full two-arrow theorem, so the displayed criterion is an explicit working formulation and the graph-specific conclusion remains an attributed claim.

The original discussion includes further goals, not conclusions of this partial realization. Lucas Sabalka asked for reasonable behavior of infinite-order elements, with block-upper-triangular or ``non-hyperbolic elements with hyperbolic mapping torus'' left as questions. Thierry Coulbois asked for naturality like Teichm\"uller space, and Martin Lustig discussed truth and beauty. Juan Souto stressed unforeseen curve-complex applications; Yair Minsky recalled applications motivating his and Masur's theorem. Proposed applications were bounded asymptotic dimension, uniformity in Dahmani--Guirardel--Osin, rigidity of lattice maps via bounded cohomology, and quasi-isometric rigidity. The notes describe the DGO statement as: for each iwip \(g\), some positive integer \(q\) makes every normal closure of \(g^{qk}\), \(k\ge1\), free and purely iwip. They attribute its graph input to Bestvina--Feighn; Vincent Guirardel questioned uniformity. No uniformity quantifier is specified here.

Martin Bridson's suggested application involving infinite-dimensional second cohomology of every infinite subgroup is retained as speculation with unspecified coefficients and ordinary-versus-bounded convention, not a theorem. Minsky suggested hyperbolic or \(\operatorname{CAT}(0)\) links and global-to-local quasi-projections for induction, and Karen Vogtmann suggested high connectivity and cohomological applications. Candidate objects listed were the free-factor complex; the trivial-edge-group splitting complex, said quasi-isometric to the sphere complex; variants of splittings over \(\Z^n\) or with rank cutoffs suggested by Matt Clay and related to length by Kasra Rafi; the poset of commensurability classes of abelian subgroups of linearly growing automorphisms, related by Souto and Bridson to Dehn multitwists and cut systems; and the Kapovich--Lustig graph. No definitions or proofs for these alternative candidates are supplied.
\subsection{Short English statement}
The aid reports the free-factor complex as a partial answer: iwips translate positively, polynomial growth gives finite vertex orbits, and connectivity plus corank-one links support some induction. Qualitative goals, alternative complexes, applications, and the even-rank co-surface obstruction remain separately qualified source claims.
\subsection{Sources}
\begin{itemize}
\item Local UnsolvedMath registry, permanent ID \(20001639\), source code \texttt{catalogue entry}; \texttt{aim-geometric-group-theory-notes.json}, \texttt{notes[147]}, tag \texttt{problem}.
\item AIM workshop ``The geometry of the outer automorphism group of a free group,'' 25--29 October 2010, edited by Johanna Mangahas, Problem 1.1: \url{https://aimath.org/WWN/outerauto/outerauto.pdf}. Source type PDF; statement status \texttt{exact}. The supplied canonical text and remarks were used; the PDF was not accessed or compiled.
\item Named inputs and discussion attributions include Bestvina--Feighn, Dahmani--Guirardel--Osin, Kapovich--Lustig, Sabalka, Bridson, Coulbois, Lustig, Souto, Minsky, Masur, Guirardel, Vogtmann, Clay, and Rafi. No authoritative full bibliographic citations or theorem identifiers for the aid's hyperbolicity, dynamics, connectivity, link, or co-surface claims are supplied. Original literature field is empty; extraction artifacts in the remarks include an editor/page header and split words.
\item Research attempt 1: \texttt{partial\_result}, status at run \texttt{partially\_solved}, rigor label \texttt{proved}, novelty confidence low; registry status \texttt{partially\_solved}, classification \texttt{PARTIAL-PROGRESS}. Created and updated 14 August 2026. Machine-generated aid, not peer reviewed; independently verify before citation.
\item Ulam AI / UnsolvedMath Contributors; CC BY 4.0.
\end{itemize}
\end{document}
